\subsection{Projective decisions over effective henselian valued fields}
\label{C:henselian-projective-subsection}

\begin{lemma}[Proved under stated hypotheses: finite projective-point decisions]
\label{C:henselian-projective-decisions}
Let $F$ be a countable henselian discretely valued field, with
decidable exact finite names, equality and arithmetic, a computable normalized
valuation $v:F^\times\to\mathbf Z$, a supplied uniformizer $\pi$,
and a supplied finite residue field with computable residue map and
lifts. Let $\widehat F$ be its completion.
There is an ordinary terminating decision for $X(F)\ne\varnothing$
for every explicitly presented smooth projective $F$-scheme $X$.
A positive answer constructs an exact point.
If $\widehat F/F$ is regular, the same assertion holds for every
explicit projective $F$-scheme, without smoothness or reducedness.
No bound for the running time or necessary congruence precision is asserted.
\end{lemma}
\begin{proof}
Write $\mathcal O=\{0\}\cup\{a:v(a)\ge0\}$.
The supplied residue lifts give complete representatives
\[
 \sum_{i=0}^{N-1}\widetilde a_i\pi^i,
 \qquad a_i\in\kappa,
 \qquad\hbox{for }\mathcal O/\pi^N\mathcal O.
                                                               \tag{C.3.HP1}
\]
To reduce an integral element, successively take its residue,
subtract its supplied lift and divide by $\pi$.
Thus each finite quotient and every polynomial congruence in it
are computable; a coefficient-field section is unnecessary.

We first verify the required point reflection.
For a smooth $F$-scheme, a point over $\widehat F$ lies in an
affine chart on which a finite list of equations has a nonzero
Jacobian minor and generates the ideal locally.
These are the usual local coordinates for smoothness.
Approximate its free coordinates and its dependent coordinates
by elements of $F$, keeping the minor and the chart denominator
nonzero. After rescaling variables and equations, all coordinates
and polynomial coefficients used in the implicit equation are integral.
Let $f$ be the square system in the dependent coordinates,
$a$ its sufficiently close approximation, $M=f'(a)$ and
$d=\det M\ne0$. For any chosen large $s$, put $\delta=d\pi^s$.
Taylor expansion gives
\[
 f(a+\delta Z)=f(a)+\delta MZ+\delta^2Q(Z),
 \qquad Q\in (Z)^2\mathcal O[Z]^c.
\]
Multiplying by $\operatorname{adj}(M)/(\delta d)$ yields
\[
 \frac{\operatorname{adj}(M)f(a)}{d^2\pi^s}
       +Z+\pi^s\operatorname{adj}(M)Q(Z)=0.             \tag{C.3.HP2}
\]
By taking $a$ sufficiently close, the constant vector is in the
maximal ideal. The Jacobian at zero is the identity.
The multivariable simple-root Hensel lemma gives a solution in
that maximal ideal. The correction $\delta Z$ can be made as
small as prescribed by $s$, so the resulting exact $F$-point
remains in the chosen chart and satisfies all its original equations.
This is the polynomial Hensel lemma for a henselian local ring,
as stated in Nowak \cite[Remark 2 and Corollary 3, p.\ 2]
{C-NowakDensity2017}; the scaling argument above gives the
nonunit-Jacobian form needed here. Hence
\[
 X(\widehat F)\ne\varnothing\quad\Longrightarrow\quad
 X(F)\ne\varnothing
                                                               \tag{C.3.HP3}
\]
for smooth $X$. If $\widehat F/F$ is regular, Jarden
\cite[Lemma 2.3 and its proof, pp.\ 6--7 of the author version]
{C-JardenHenselian1994} proves that $F$ is existentially closed in
$\widehat F$. Apply that statement to an affine chart of a
projective point, including the inverse of its nonzero chart
coordinate. It proves (C.3.HP3) for arbitrary projective schemes.

Embed the given scheme in $\mathbf P_F^m$ using its supplied
homogeneous equations. Multiply each equation by a power of
$\pi$ to make its coefficients integral.
In the positive search enumerate all exact vectors of $F^{m+1}$
with one designated coordinate equal to one, and test every
equation. These charts enumerate all projective $F$-points.

In the negative search, for $N=1,2,\ldots$, exhaust the finite
set of vectors in $(\mathcal O/\pi^N)^{m+1}$ having at least
one unit coordinate, and retain those satisfying every equation.
An empty level proves that no point exists even over $\widehat F$.
Conversely, if every level is nonempty, reduction makes these
finite sets into a finitely branching tree with vertices at every
depth. An infinite branch gives a compatible primitive vector
in $\widehat{\mathcal O}^{m+1}$. Polynomial continuity and
separatedness show that it satisfies every equation, and its
unit coordinate makes it a projective point over $\widehat F$.
By (C.3.HP3) an exact $F$-point then exists.

Dovetail the two searches in finite steps. If a point exists
the positive search terminates. If none exists, point reflection
and the preceding compactness argument force an empty finite
congruence level, so the negative search terminates.
Both answers are correct and one necessarily occurs.
The zero-dimensional and empty-scheme cases use the same argument.
\end{proof}

\begin{theorem}[Proved under stated hypotheses: actual central division factors]
\label{C:henselian-central-simple-algorithms}
Let $F$ have the effective valued-field data of
Lemma~\ref{C:henselian-projective-decisions}.
For every explicitly supplied finite central simple $F$-algebra
$A$, there are ordinary algorithms computing its index and
an actual isomorphism $A\simeq M_r(D)$, including the
multiplication table of the central division algebra $D$,
all matrix units and primitive idempotents.
There is no prime-to-characteristic restriction.
If factorization over all finite towers over $F$ is also supplied,
reduced norms and reduced characteristic polynomials are computable.
\end{theorem}
\begin{proof}
Put $\dim_F A=n^2$. For $1\le j\le n$, form the closed
subscheme $X_j$ of $\operatorname{Gr}_{nj}(A)$ consisting
of the $nj$-dimensional left ideals.
This is an explicit projective scheme. In Pluecker coordinates
$w\in\bigwedge^{nj}A$ use the usual Grassmannian equations
and, for every basis element $b$ of $A$ and every basis
$(nj-1)$-covector $\xi$, impose
\[
             L_b(\iota_\xi w)\wedge w=0.             \tag{C.3.HP4}
\]
For a decomposable nonzero $w$ the contractions span its
subspace, so these equations express precisely left stability.

The scheme $X_j$ is smooth. Indeed after a finite separable
splitting extension it is the ordinary Grassmannian
$\operatorname{Gr}_j(F^n)$: for $A=M_n$, a stable direct
summand is $F^n\otimes W$, with $W$ a rank-$j$ direct summand
of the row module. Matrix idempotents and matrix units give
this correspondence over every algebra over the splitting field,
so this is an identification of schemes, not only of field points.
Grassmannian graph charts are affine spaces; smoothness descends
along the splitting field extension.
Apply Lemma~\ref{C:henselian-projective-decisions} to every $X_j$.

If $A\simeq M_r(D)$ with $\dim_FD=d^2$ and $n=rd$,
every left ideal has $F$-dimension a multiple of $nd$,
and all multiples $nd,2nd,\ldots,n^2$ occur.
This follows directly by viewing left ideals as sums of
simple column modules, each of dimension $rd^2=nd$.
Consequently
\[
       d=\min\{j\ge1:X_j(F)\ne\varnothing\}.          \tag{C.3.HP5}
\]
This computes the index without a period--index theorem.

Enumerate $e\in A$ until $e^2=e$ and
$\dim_F Ae=nd$; both tests are finite linear algebra.
A primitive matrix idempotent supplies such an $e$, so the
search terminates. Then $D=eAe$ is a division algebra, with
its actual multiplication table read from a computed basis.
Put $r=n/d$. Enumerate finite lists $u_{ij}\in A$ with
$u_{11}=e$, $u_{ij}u_{kl}=\delta_{jk}u_{il}$ and
$\sum_i u_{ii}=1$. Such matrix units exist by the matrix
decomposition. They give the mutually inverse maps
\[
 A\longrightarrow M_r(D),\quad
 a\longmapsto(u_{1i}au_{j1})_{ij},\qquad
 (d_{ij})\longmapsto\sum_{i,j}u_{i1}d_{ij}u_{1j}.
\]
All searches use exact arithmetic and a proved existing witness.

For reduced norms, enumerate finite separable tower fields and
actual matrix-algebra splitting isomorphisms. A finite separable
splitting field exists for every central simple algebra
\cite[Proposition 11.8.5, Tag 0752]{C-StacksGlobalFunctionArithmetic}.
The determinant and characteristic polynomial in such a splitting
lie in $F$ by invariance under inner conjugation and separable
descent. Their ground-field coordinates are read from the tower
basis, exactly as in
Theorem~\ref{C:global-function-central-simple-algorithms}.
The supplied factorization verifies all field candidates.
This proves the final assertion, including wild degrees.
\end{proof}

\begin{theorem}[Proved under stated hypotheses: all hermitian isometry decisions]
\label{C:henselian-hermitian-isometry}
Let $F$ have the effective data of
Lemma~\ref{C:henselian-projective-decisions}, with
$\operatorname{char}F\ne2$ and $\widehat F/F$ regular.
Let $D$ be a finite-dimensional division algebra, supplied by
a multiplication table over its center $L$, and let $\sigma$
be an involution with $L^\sigma=F$, supplied as an $F$-linear map.
Thus $L=F$ or $L/F$ is a separable quadratic extension.
For each $\varepsilon\in\{1,-1\}$, isometry of nonsingular
$\varepsilon$-hermitian forms over $(D,\sigma)$ is decidable
by an ordinary algorithm. Positive decisions construct actual
isometry matrices.

Isometry classes have a decidable literal least-Gram-code
index set. Their orthogonal-sum Grothendieck group has
decidable literal least-pair coordinates, computable group
operations and actual positive/negative form representatives.
No local classification or unknown hermitian-isometry quotient
is part of the output.
\end{theorem}
\begin{proof}
First compare the $D$-ranks; unequal ranks are not isometric.
For common rank zero return the unique isometry of zero spaces.
For common rank $n>0$ put $W=D^{2n}$ and
$b=h\perp(-g)$, a nonsingular form on $W$.
Let $a=\dim_FD$. In the ordinary $F$-Grassmannian
$\operatorname{Gr}_{na}(W)$ impose two closed conditions:
the universal subspace is stable under right multiplication
by every element of a fixed $F$-basis of $D$, and its
restriction of $b$ is identically zero.
The first condition uses the contraction equations
(C.3.HP4), now for right multiplications.
For the second, take every pair of basis
$(na-1)$-covectors $\xi,\eta$ and require every $F$-coordinate
of
\[
                     b(\iota_\xi w,\iota_\eta w)=0.   \tag{C.3.HP6}
\]
The form $b$ is $F$-bilinear because $F$ is fixed by $\sigma$.
These are explicit homogeneous quadratic equations in the
Pluecker coordinates. The resulting scheme $Y_b$ is projective.
Its $F$-points are exactly the right $D$-subspaces of $W$
of $D$-dimension $n$ on which $b$ vanishes: a finite
right module over the division algebra is free, and its
$F$-dimension is its $D$-rank times $a$.
We use the arbitrary-projective clause of
Lemma~\ref{C:henselian-projective-decisions}; in particular no
assumption about geometric ranks after a unitary-center
splitting, or smoothness of this presentation, is necessary.

We recall why
\[
       Y_b(F)\ne\varnothing
         \quad\Longleftrightarrow\quad h\simeq g.     \tag{C.3.HP7}
\]
A totally isotropic half-dimensional subspace has an
isotropic complement after adjusting any dual complement:
the Gram matrix initially has shape
$\left(\begin{smallmatrix}0&I\\\varepsilon I&C\end{smallmatrix}\right)$
with $C^*=\varepsilon C$. Replacing the second block of
basis vectors by itself plus the first block times
$-\varepsilon C/2$ makes $C$ zero.
Thus a metabolic nonsingular form over $D$ is hyperbolic.
Conversely a hyperbolic form has its standard half-dimensional
isotropic subspace.
Finally Witt cancellation for hermitian forms over a division
algebra with $2\ne0$, with the exact hypotheses already used
in the checked QSS transfer theorem
\cite[Section 3, pp.\ 271--273]{C-QSS1979}, identifies equal
rank and hyperbolicity of $h\perp(-g)$ with actual isometry
of $h$ and $g$.
This also treats rank zero, and alternating forms when they
occur. It proves (C.3.HP7).

The projective-point decision now decides isometry.
If it is positive, enumerate matrices $P,Q$ over $D$ and
test $PQ=QP=I$ and $P^*HP=G$; the proved isometry guarantees
termination and supplies an actual matrix.

Encode all valid nonsingular Gram matrices by natural numbers
using the exact field and algebra names. Validity, rank and
isometry are decidable. Retain a code exactly when no smaller
valid code is isometric to it. This is a finite predecessor
test and gives a decidable repetition-free index set.
Witt cancellation makes the orthogonal-sum monoid cancellative.
Two pairs $(H_+,H_-)$ and $(H'_+,H'_-)$ give the same group
element exactly when
\[
                  H_+\perp H'_-\simeq H'_+\perp H_- . \tag{C.3.HP8}
\]
We have just proved this test decidable. Retain the least
valid pair in its equivalence class; every pair has only
finitely many predecessors, so this gives literal decidable
group coordinates. Block sum and interchange compute addition
and inverse, followed by the finite canonicalization test.
The pair itself is an actual positive/negative representative.
All conclusions are ordinary terminating algorithms with the
stated effective data, not decisions from arbitrary streams.
\end{proof}
